% =====================================================================
%  Thème 3 — CORRIGÉ — Niveau MOYEN
% =====================================================================
\documentclass[11pt,a4paper]{article}
\usepackage[utf8]{inputenc}
\usepackage[T1]{fontenc}
\usepackage{lmodern}
\usepackage[french]{babel}
\usepackage{amsmath,amssymb,mathtools}
\usepackage{icomma}
\usepackage{xcolor}
\usepackage[most]{tcolorbox}
\usepackage{enumitem}
\usepackage{array}
\usepackage{booktabs}
\usepackage{fancyhdr}
\usepackage[hidelinks]{hyperref}
\usepackage[a4paper,margin=2.2cm,headheight=26pt,headsep=10pt]{geometry}

\definecolor{primary}{RGB}{223,87,99}
\definecolor{primarydark}{RGB}{150,42,55}
\definecolor{excol}{RGB}{0,135,90}

\tcbset{boxbase/.style={enhanced,breakable,boxrule=0.9pt,arc=2.5pt,
   left=3mm,right=3mm,top=2.3mm,bottom=2.3mm,
   fonttitle=\bfseries,coltitle=white,toptitle=0.6mm,bottomtitle=0.6mm}}
\newtcolorbox[auto counter]{correction}[1][]{boxbase,colback=excol!5,colframe=excol,
  title={\textsc{Corrigé}~\thetcbcounter\ifx\relax#1\relax\relax\else\textmd{ --- \itshape #1}\fi}}

\pagestyle{fancy}
\fancyhf{}
\fancyhead[L]{\small\textcolor{primarydark}{\textbf{Fiche 6} — Les limites}}
\fancyhead[R]{\small\textcolor{primarydark}{\textsc{Thème 3 — Corrigé Moyen}}}
\fancyfoot[C]{\small\thepage}
\renewcommand{\headrulewidth}{0.8pt}
\renewcommand{\headrule}{\hbox to\headwidth{\color{primary}\leaders\hrule height \headrulewidth\hfill}}

\newcommand{\R}{\mathbb{R}}

\begin{document}

\begin{center}
{\Large\bfseries\color{primarydark}Thème 3 — Limites à l'infini : le terme dominant}\\[2pt]
{\large\color{excol}Corrigé — Niveau Moyen}
\end{center}
\vspace{6pt}

\begin{correction}[le même polynôme des deux côtés]
$P(x)=-2x^3+4x^2-1$, terme dominant $-2x^3$.
\begin{enumerate}[label=\alph*),leftmargin=2em,itemsep=3pt]
  \item $\displaystyle\lim_{x\to+\infty} P(x)=\lim_{x\to+\infty}(-2x^3)=-2\times(+\infty)=-\infty$.
  \item $\displaystyle\lim_{x\to-\infty} P(x)=\lim_{x\to-\infty}(-2x^3)=-2\times(-\infty)=+\infty$.
  \item Oui, les résultats sont \emph{opposés} ($-\infty$ et $+\infty$) : le degré $3$ est \textbf{impair},
        donc $x^3$ change de signe entre $-\infty$ et $+\infty$.
\end{enumerate}
\end{correction}

\begin{correction}[rationnelles : les trois cas]
\begin{enumerate}[label=\alph*),leftmargin=2em,itemsep=3pt]
  \item $n=m=2$ : $\displaystyle\lim_{x\to+\infty}\frac{6x^2-x+2}{2x^2+5}=\frac{6x^2}{2x^2}=\frac{6}{2}=3$.
  \item $n=2<m=3$ : $\displaystyle\lim_{x\to-\infty}\frac{x^2+3}{x^3-x}=\lim_{x\to-\infty}\frac{x^2}{x^3}=\lim_{x\to-\infty}\frac{1}{x}=0$.
  \item $n=3>m=2$ : $\displaystyle\lim_{x\to-\infty}\frac{4x^3-x}{2x^2+1}=\lim_{x\to-\infty}\frac{4x^3}{2x^2}=\lim_{x\to-\infty}2x=-\infty$.
\end{enumerate}
\end{correction}

\begin{correction}[asymptote horizontale]
$f(x)=\dfrac{5x^2-3}{2x^2+x+1}$, degrés égaux ($n=m=2$).
\begin{enumerate}[label=\alph*),leftmargin=2em,itemsep=3pt]
  \item $\displaystyle\lim_{x\to+\infty} f(x)=\lim_{x\to+\infty}\frac{5x^2}{2x^2}=\frac{5}{2}$, et de même
        $\displaystyle\lim_{x\to-\infty} f(x)=\frac{5}{2}$.
  \item La courbe se rapproche de la droite $y=\dfrac{5}{2}$ aux deux bouts : \textbf{asymptote horizontale
        $y=\tfrac52$}.
\end{enumerate}
\end{correction}

\begin{correction}[justification par mise en facteur]
\begin{enumerate}[label=\alph*),leftmargin=2em,itemsep=3pt]
  \item $\displaystyle\frac{3x^2+x}{x^2-4}=\frac{x^2\bigl(3+\frac1x\bigr)}{x^2\bigl(1-\frac{4}{x^2}\bigr)}$.
  \item Pour $x\neq 0$, $\dfrac{x^2}{x^2}=1$, et quand $x\to+\infty$ on a $\dfrac1x\to 0$, $\dfrac{4}{x^2}\to 0$ :
        \[ \lim_{x\to+\infty}\frac{3+\frac1x}{1-\frac{4}{x^2}}=\frac{3+0}{1-0}=3. \]
        On retrouve bien $3$, le rapport des coefficients dominants.
\end{enumerate}
\end{correction}

\begin{correction}[attention au signe du côté $-\infty$]
\begin{enumerate}[label=\alph*),leftmargin=2em,itemsep=3pt]
  \item $\displaystyle\lim_{x\to-\infty}\frac{x^4-1}{-2x^2}=\lim_{x\to-\infty}\frac{x^4}{-2x^2}=\lim_{x\to-\infty}\frac{x^2}{-2}$.
        Or $x^2\to+\infty$, donc la limite vaut $\dfrac{+\infty}{-2}=-\infty$.
  \item $\displaystyle\lim_{x\to-\infty}\frac{3x^5}{x^4+1}=\lim_{x\to-\infty}\frac{3x^5}{x^4}=\lim_{x\to-\infty}3x=-\infty$.
\end{enumerate}
\end{correction}

\end{document}
